\chapter{Conclusions}
In this laboratory work, I implemented the Romanian Playfair cipher using
a keyword matrix with six rows and five columns. Merging J into I made
room for the 31-letter Romanian alphabet while keeping Î as a separate
letter. The program handles uppercase and lowercase input, Romanian
cedilla variants and ordinary message spaces.

I verified the three transformation rules and their inverses: horizontal
shifts, vertical shifts and rectangle substitution. The examples also
confirmed the importance of preparing pairs correctly. When two current
letters are identical, only the first original letter is consumed; the
second must be reconsidered. X is the usual filler, with Q used when an
X would otherwise create an identical pair.

The console output makes the process easy to check: it shows the matrix,
prepared pairs and one short explanation for each transformation.
Decryption correctly recovers the prepared plaintext. Restoring word
spaces, distinguishing an original J from I and resolving filler letters
remain manual tasks because this information is not preserved by the
cipher. The work reinforced my understanding of pair substitution, modular
wraparound and input validation through a small, readable implementation.
